\section{Bounds for the polynomial cubic adjacent pencil}

\begin{finding}{P-18}{a twenty-seven-term border upper bound and exact lower certificates for the cubic polynomial seed}
Let $M_-$ and $M_+$ be the first-factor-labelled ordered
coefficient tensors of the two $3\times3$ determinants,
respectively permanents, in the columns $012$ and $013$
of a $3\times4$ matrix. Their format is $2\times12\times12\times12$.
Over every algebraically closed field,
\[
 \BR(M_-)\leq27,\qquad 21\leq\BR(M_+)\leq24.
\]
If the characteristic is not three, also $\BR(M_-)\geq21$.
In characteristic two the two tensors coincide, so their
common border rank is between $21$ and $24$.
In particular, in characteristic zero the bounds are
$21\leq\BR(M_-)\leq27$ and $21\leq\BR(M_+)\leq24$.
No exact value for either polynomial pencil is asserted.
\end{finding}
\begin{proof}
Write $V=k^3$ and let $E$ have basis $e_{r,c}$ with
$0\leq r\leq2$ and $0\leq c\leq3$. Let
$X,Y,Z_0,Z_1:V\to E$ be the split embeddings into
columns $0,1,2,3$. Define the linear map
\[
 \Phi:V^{\otimes3}\longrightarrow k^2\otimes E^{\otimes3}
\]
by
\[
 \Phi(u\otimes v\otimes w)
 =\sum_{i=0}^1a_i\otimes
   \operatorname{per}\bigl(X(u),Y(v),Z_i(w)\bigr),
\]
where $\operatorname{per}$ means the six ordered tensor
words, including their multiplicities. This expression
is trilinear in $u,v,w$, hence defines the indicated
linear map. Expanding in the row basis gives
\[
 \Phi(D_3)=M_-,\qquad \Phi(P_3)=M_+.
\]
This is an identity of integral coefficient tensors;
no division by $3!$ or polynomial polarization is used.

For a pure input tensor, P-10 with $b=3,j=2$ supplies
a six-term border decomposition of its image. Therefore
\[
 \BR(\Phi(T))\leq6\BR(T)
 \tag{P18.1}
\]
for every $T$: the inverse image under $\Phi$ of the
closed $6r$-secant cone contains every sum of $r$
pure tensors and hence its closure. The one-label
specialization of P-11's determinant tangent formula
has five terms, giving $\BR(D_3)\leq5$ over
$\mathbb Z$ after parameter cancellation. P-10 gives
$\BR(P_3)\leq4$ in every characteristic. Equation
(P18.1) proves auxiliary upper bounds $30$ and $24$.

We improve the determinant upper bound to $27$ by a
partially symmetric intermediate space. First, if
$\operatorname{char}k\ne2$, every element of
$F\otimes\Gamma^2H$, where $\dim F=3$ and $\dim H=6$,
is a limit of sums of nine tensors $f\otimes v^{\otimes2}$.
Here $\Gamma^2H=(H^{\otimes2})^{\mathfrak S_2}$ is the
space of invariant tensors; its coordinates are the
coefficients at $(i,j)$ for $0\leq i\leq j\leq5$.
To prove this assertion consider the polynomial map
\[
 \Theta:(F\times H)^9\longrightarrow F\otimes\Gamma^2H,
 \qquad ((c_h,v_h))\longmapsto\sum_{h=0}^8c_h\otimes v_h^{\otimes2}.
 \tag{P18.3}
\]
Let $f_0,f_1,f_2$ denote the standard basis of $F$,
and use the following explicit point of the source:
\[
\begin{split}
 &c_h=f_{\lfloor h/3\rfloor}\quad(0\leq h\leq8),
 \qquad v_h=e_h\quad(0\leq h\leq5),\\
 &v_6=(1,1,0,1,1,0),\quad
 v_7=(1,0,1,2,0,1),\quad
 v_8=(0,1,1,0,2,2).
\end{split}
\]
Thus the nine labels are three copies of each basis
vector. Order the target
coordinates by $(a,i,j)$ with $a$ first and $i\leq j$
lexicographically; order the source coordinates by
the summand $h$, its three label coordinates, and
then its six vector coordinates. The Jacobian entries
are, explicitly,
\[
 \frac{\partial\Theta_{aij}}{\partial c_{hb}}
 =\delta_{ab}v_{hi}v_{hj},\qquad
 \frac{\partial\Theta_{aij}}{\partial v_{h\ell}}
 =c_{ha}(\delta_{i\ell}v_{hj}+\delta_{j\ell}v_{hi}).
\]
The certificate \path{nine_square_jacobian_certificate.json}
specifies $63$ columns of this $63\times81$ matrix;
their exact integer determinant is $-16$.
The standard-library checker
\path{verify_nine_square_jacobian.py} reconstructs the
point and every Jacobian entry from these formulas
and computes that determinant by fraction-free elimination.

For completeness, this Jacobian certificate proves
dominance without a characteristic-zero assumption.
Fix the other eighteen source coordinates at the
displayed point. Translate that point and its image
to zero. The resulting polynomial map from $\mathbb A^{63}$
to itself has an invertible linear term whenever
$\operatorname{char}k\ne2$. Solving one homogeneous
degree at a time gives a formal power series inverse.
A polynomial relation vanishing on its image would,
after composition with that inverse, be the zero
formal polynomial, and hence the zero polynomial.
Thus (P18.3) has dense image in these characteristics,
which proves the nine-term border assertion.

Apply this fact with $H=X(V)\oplus Y(V)$, the common
two columns of the two determinants. In the cofactor
expansion
\[
 \det[X,Y,Z_i]=\sum_{r=0}^2 (Z_i)_r\,q_r(X,Y),
\]
let $Q_r\in\Gamma^2H$ be the ordered coefficient tensor
of the signed quadratic cofactor $q_r$. In particular,
an off-diagonal monomial with coefficient $c$ contributes
$c(u\otimes v+v\otimes u)$; no coefficient is divided
by two. Put $Q=\sum_r f_r\otimes Q_r$ and identify
$Z_i(f_r)=e_{r,2+i}$. Define the linear map
\[
 \Psi:F\otimes\Gamma^2H\longrightarrow k^2\otimes E^{\otimes3}
\]
by inserting $Z_i(f)$ in each of the three positions
beside the two ordered factors of a quadratic tensor,
and summing with label $a_i$. In particular,
\[
 \Psi(f\otimes v^{\otimes2})=
 \sum_{i=0}^1a_i\otimes
 \bigl(v\otimes v\otimes Z_i(f)+v\otimes Z_i(f)\otimes v
                 +Z_i(f)\otimes v\otimes v\bigr).
\]
The cofactor expansion gives $\Psi(Q)=M_-$. Each
displayed image has border rank at most three: it
is the limit at zero of
\[
 \frac1t\left[\sum_{i=0}^1a_i\otimes
       (v+tZ_i(f))^{\otimes3}
       -(a_0+a_1)\otimes v^{\otimes3}\right].
 \tag{P18.4}
\]
There are three pure summands for $t\ne0$; their sum
is regular at zero. Consequently the inverse image
under $\Psi$ of the closed border-rank-$27$ cone
contains every sum of nine labelled squares and
therefore their entire closure. This proves
$\BR(M_-)\leq27$ in characteristic different from two.
In characteristic two $M_-=M_+$, and the already
proved permanent upper bound $24$ is stronger. Thus
the asserted determinant upper bound holds in every
characteristic.

There are also explicit one-parameter certificates
for the auxiliary determinant bound $30$ and the
permanent bound $24$.
The five-term row determinant curve has scalar poles
of order at most one. Apply the six-point difference
to each pure row term, with its difference parameter
$s=t^2$. After the difference is divided by $s^2$,
its error beyond the desired ordered cubic is divisible
by $s$. Multiplication by a row coefficient with a pole
of order at most one leaves an error divisible by $t$.
Thus the resulting thirty-term expression is regular
at zero with limit $M_-$. For the four-term row permanent
curve the scalar poles have order at most two, and
$s=t^3$ gives the corresponding twenty-four-term curve.
The standalone program
\path{verify_polynomial_cubic_upper.py} reconstructs both
curves and checks all $2\cdot12^3=3456$ constant
coefficients and every negative-power cancellation over
$\mathbb Z$. Its explicit atom lists are recorded in
\path{polynomial_cubic_upper_certificate.json}.

For the lower bounds use the double-Koszul map
\[
 K_M:A\otimes\Lambda^4E\otimes E^*
 \longrightarrow\Lambda^2A\otimes\Lambda^5E\otimes E,
 \qquad \dim A=2,
\]
given, for $M=\sum m_{a b c d}a_a\otimes e_b\otimes e_c\otimes e_d$, by
\[
 \alpha\otimes\eta\otimes\beta
 \longmapsto\sum m_{a b c d}
 (a_a\wedge\alpha)\otimes(e_b\wedge\eta)
                       \otimes\beta(e_c)e_d.
\]
Its matrix has $9504$ rows and $11880$ columns. Its rank
on a nonzero pure tensor is
$\binom{11}{4}=330$. Consequently
\[
 \BR(M)\geq\left\lceil\frac{\rank K_M}{330}\right\rceil.
 \tag{P18.2}
\]

The following table summarizes exact integer minors
in $905$ different weight blocks for each tensor,
with $1810$ certified minors in total:
\[
\begin{array}{c|r|l}
\text{tensor and certificate}&\text{sum of minor orders}
 &\text{nonvanishing condition}\\\hline
M_-&6680&\text{each determinant is a unit times }2^u3^v\\
M_+&6626&\text{every determinant is a unit}
\end{array}
\]
Here $u,v\geq0$, and a unit means $1$ or $-1$.
To explain why the block orders add, give $e_{r,c}$
weight $R_r+C_c$ and $a_i$ weight
\[
 -R_0-R_1-R_2-C_0-C_1-C_{2+i}.
\]
Every coefficient term of $M$ has total weight zero.
The map $K_M$ is equivariant. Its source weight at
$a_\alpha\otimes e_I\otimes e_k^*$ is
$\operatorname{wt}(a_\alpha)+\sum_{i\in I}\operatorname{wt}(e_i)
-\operatorname{wt}(e_k)$; its target weight at
$(a_0\wedge a_1)\otimes e_J\otimes e_l$ is
$\operatorname{wt}(a_0)+\operatorname{wt}(a_1)
+\sum_{j\in J}\operatorname{wt}(e_j)+\operatorname{wt}(e_l)$.
Different weights therefore contribute disjoint matrix
blocks, in every characteristic.

The complete row and column indices and integer
determinants are in
\path{polynomial_cubic_lower_certificate.json}.
The independent standard-library verifier
\path{verify_polynomial_cubic_lower.py} reconstructs
every selected entry from the displayed map, verifies
the weights and their distinctness, and evaluates all
minors by exact fraction-free elimination. It also
checks the prime divisors stated in the table. Thus
the first row gives rank at least $6680$ when the
characteristic is neither two nor three; the second
gives $6626$ in every characteristic. Both exceed
$20\cdot330=6600$, so (P18.2) gives the claimed lower
bounds. In characteristic two $M_-=M_+$, which also
covers the determinant there. The displayed minors
do not establish the determinant lower bound twenty-one
in characteristic three, and no such assertion is made.
\end{proof}

\paragraph{What this leaves open.}
P-02 identifies these tensors with the polynomial
$b=3,j=2$ minima for every $n\geq4$. At $n=4$ the
flattening image has dimension $16$. The permanent
upper bound therefore caps the result obtainable from
this particular substitution at $16-2+24=38$.
The determinant upper bound caps it at
$16-2+27=41$. Thus this particular two-plane cubic
substitution cannot improve the proved determinant
lower bound $43$, regardless of the unresolved exact
value of the cubic pencil. The present lower
certificates give $21$ for that pencil in characteristic
different from three.

\paragraph{Geometric dictionary for P-18.}
The maps $\Phi$ and $\Psi$ are maps of free integral
modules and their associated tensor bundles. The
Jacobian minor proves density for a map of affine
schemes after inverting two. The upper curves
are integral Segre sections on $\mathbb G_m$ whose
grouped sum extends over $\mathbb A^1$. The Koszul
map is a morphism involving exterior powers of locally
free sheaves. The split-torus characters give its
weight-block decomposition, and the actual integer
minors carry its rank information. These are the
EGA~I \S9.5 image/closure construction, the SGA~3
Expos\'e~I \S\S4.5--4.6 module/group schemes, and
the SGA~6 Expos\'e~V Definition~2.1, Examples~2.2(b)
and Lemma~2.2.1 exterior operations, with the exact
printed and facsimile locators in the common dictionary.
Neither a $K_0$ class nor an unproved smoothability
claim replaces these coefficient and minor checks.
